\section{Related work}
\label{related-work}

% \subsection{Behavioral and psychometric evaluation of LLMs}
The NLP community has adopted frameworks from the behavioral sciences to study LLMs' latent traits \citep{ward-etal-2024-evaluating, perez-etal-2023-discovering,wadi2026does}. 
Recent works have utilized psychometric tests to measure a wide array of constructs \citep{jiang-etal-2024-personallm, li-etal-2024-evaluating-psychological, helwe-etal-2025-navigating}. 
However, recent studies demonstrate that eliciting psychometric responses via text generation is highly unstable. 
Simple prompt perturbations can significantly downgrade the consistency of LLM responses on persona axes \citep{shu-etal-2024-dont}. 
Moreover, traditional psychometric testing on LLMs suffers from low ecological validity when relying on standard generative outputs \citep{jung-etal-2026-psychometric}. 
Part of this instability stems from generative sampling, which is a methodological artifact borrowed from human surveying,
and obscures their actual underlying distribution \citep{deas-mckeown-2025-artificial}. 

\paragraph{Uncertainty and probability in LLMs}
Recognizing the limitations of text-based sampling, a growing body of literature has turned to analyzing exact token probabilities to better understand LLMs' internal states. 
\citet{kuribayashi-etal-2024-psychometric} demonstrated that utilizing exact next-word probabilities is far superior to instruction-tuned prompting when simulating human cognitive behaviors. 
Similarly, recent work cautions that prompting LLMs for open-ended responses yields poorly calibrated text \citep{liu2025badworktimecrosscultural}, 
whereas their exact internal token probabilities remain highly predictive 
\citep{bentegeac2025token, kadavath2022language}.
To quantify the variance within these predictive probability distributions, 
the NLP community relies almost exclusively on Shannon entropy or its semantic derivatives \citep{farquhar2024detecting}. 
Entropy is highly effective for tasks where the output space consists of distinct, categorical facts (e.g., detecting hallucinations across mutually exclusive answers). 
However, standard metrics like entropy fail when applied to behavioral evaluation using psychometric rating scales. 
This is because entropy treats all tokens as unordered, categorical variables, ignoring the ordinality of a Likert scale. 
Because standard uncertainty metrics are indifferent to this ordinality, we lack the mathematical foundation to utilize exact token probability mass functions (PMFs) for behavioral evaluation. 
As demonstrated in \ref{sec:methods}, our framework resolved this impasse by formulating a measure of consensus specifically designed for ordinal variables.

\paragraph{Geopolitical bias and ethnocentrism in LLMs}

% A critical application for our robust behavioral framework is the measurement of latent geopolitical and cultural biases. 
As LLMs are deployed globally, there is growing concern regarding their cultural neutrality. 
A robust body of literature documents that LLMs exhibit systemic 
geopolitical and geographical biases \citep{shwartz-2022-good, bhatia-shwartz-2023-gd, bhatia-etal-2024-local, naous-xu-2025-origin, kruspe-2024-musical}.
Furthermore, models frequently struggle to bridge culture gaps \citep{tanwar-etal-2025-know, liu-etal-2024-multilingual, saha-etal-2025-reading, liu-etal-2026-tailored}.
To systematically quantify these cultural alignments, researchers frequently turn to global sociological instruments, such as the World Values Survey \citep{zhao-etal-2024-worldvaluesbench}. Yet, these evaluations suffer from the exact generative brittleness highlighted above. 
\citet{kabir-etal-2025-break} noted that evaluating cultural alignment in LLMs using standard multiple-choice surveys leads to wildly inconsistent outputs under minor perturbations, such as choice reordering. 

To bypass this brittleness, research has advocated moving toward unconstrained, open-ended generative evaluations. 
However, relying on open-ended generation introduces a new set of methodological bottlenecks. 
Specifically, it is notoriously difficult to reliably parse distinct categorical or ordinal scores from free-text outputs. 
Open-ended generation is highly susceptible to "fence-sitting", where models generate evasive disclaimers (e.g., "As an AI, I do not have a country...") rather than revealing their latent behavioral stance \citep{bhatia2025value, liu2025badworktimecrosscultural}. 
Dealing with these nuances typically requires utilizing another LLM as a judge to evaluate the text, which merely injects a secondary layer of bias into the evaluation pipeline.


